\chapter[Genomes Copy, Borrow, and Remodel]{Genomes Copy,\protect\\ Borrow, and Remodel}

\begin{kaichang}
As autumn turns to winter, roasted-sweet-potato ovens return to street corners. Break open a syrupy roasted sweet potato and take a bite: sweet and fluffy. You may not know that, strictly speaking, it is a \textbf{naturally transgenic organism}.

In 2015, scientists examined nearly three hundred cultivated sweet-potato varieties from around the world. In \textbf{every one}, they found genomic fragments left by the same soil bacterium, Agrobacterium. Long ago, this bacterium ``stuffed'' DNA into cells of a sweet-potato ancestor. Rather than causing trouble, the foreign genes were retained and passed down into every roasted sweet potato today.

No laboratory, scientists, or pipettes were involved. Nature had performed ``genetic modification'' itself. What sort of machine, then, is a genome? Is it not supposed to be a carefully protected blueprint copied letter for letter? This chapter explores the ``book of life'' as it really is: whole pages get photocopied, passages move, pages are borrowed from neighbors, and old parts are hammered into entirely new uses.
\end{kaichang}

\section{A Restless Book}

So far, you may have acquired this picture: DNA is a precise instruction manual (Chapter 70), copying involves proofreading scribes (Chapter 72), and mutations are mostly spelling errors---substituted or missing letters. All this is correct, but it describes only the smallest changes, at the single-letter level.

Focusing on letters alone misses the book's more dramatic side. Real genomes constantly undergo large-scale changes:

\begin{itemize}[itemsep=2pt]
\item Entire passages are \textbf{photocopied} into two, three, or even dozens of copies.
\item Some paragraphs \textbf{move themselves}, jumping from one chapter to another.
\item Sometimes a page is simply \textbf{torn from another book} and pasted in.
\item \textbf{Scraps} left by viruses that ``died'' hundreds of millions of years ago remain tucked inside, and some are even put back to work.
\end{itemize}

This sounds like a librarian's nightmare. But recall Chapter 76: variation is evolution's raw material. A misspelled letter is variation; copying, moving, and borrowing create variation on a larger scale. How do these major operations happen? What consequences do they usually have? Which happened to become key evolutionary inventions?

One warning first: genomes do \textbf{not} make these changes ``in order to'' evolve. No molecule knows what the future will need. Changes simply keep occurring through chemistry and probability, and environment and time do the filtering. That causal sequence must never be reversed.

\section{Photocopying a Passage: Gene Duplication}

Begin with a fact you can verify with your own eyes, or at least in a mirror: you can distinguish red from green. Your dog probably cannot.

A dog's world is roughly mixed from two colors, whereas you have three ``color sensors'': retinal opsins most sensitive to blue, green, and red. Genes for red and green opsins sit side by side on the X chromosome and have about 98\% sequence similarity. Why so alike? They began as \textbf{two copies of one gene}.

Some thirty or forty million years ago, our primate ancestors experienced a ``photocopying accident.'' During sperm or egg production, chromosomes pair and exchange segments through meiosis, as Chapter 69 described. Sometimes they misalign: gene A on one chromosome lines up with gene B on the other. Exchange then gives one cell an extra segment and deprives the other of one. Such ``unequal crossing over'' duplicates an entire stretch of DNA. Red and green opsin genes are twins from one ancient duplication. They subsequently accumulated separate mutations: one became sensitive to red light, the other remained green-sensitive, adding a color dimension to our ancestors' world. Among mammals with only two opsins, distinguishing ripe red fruit from green leaves provided a tangible survival advantage.

This is \textbf{gene duplication}: an entire DNA segment, sometimes one gene and sometimes several, acquires extra copies. It is one of evolution's most important raw-material production lines for a simple reason. A lone original is busy doing its job, and a sufficiently large mutation may spoil it and be immediately removed by selection. With two copies, one can continue the old job while the other becomes ``free'' to accumulate mutations. Usually it becomes a useless defective copy, a pseudogene; occasionally, it stumbles upon a useful new function.

A whole family in your body arose through this copying and divergence. Hemoglobin from Chapter 54 comes in several versions. Fetal hemoglobin differs from adult hemoglobin and is better at ``taking'' oxygen from maternal blood. The $\alpha$- and $\beta$-globin gene families governing these versions arose through repeated copying and remodeling over billions of years. Myoglobin, which stores oxygen in muscles, is an older member of the duplication chain.

Duplication can even leave traces of what you eat. Salivary amylase, Chapter 56's enzyme that makes chewed starch taste sweet, is produced by \emph{AMY1}. People differ greatly in copy number: some have only two copies, others more than ten. Studies find higher average copy numbers and more salivary amylase in populations with long histories of starch-rich diets, including rice, wheat, and tubers. Even the cultural history of ``what mainly fills your bowl'' can leave a genomic record of photocopy counts.

\begin{qiaoliang}
How much of your genome consists of ``photocopies'' and ``immigrants''? Calculate using roughly $3\times 10^{9}$ base pairs in a human haploid genome:

\begin{itemize}[itemsep=2pt]
\item Protein-coding sequence occupies only about $1.5\%$, roughly 45 million bases.
\item ``Mobile passages,'' the transposable elements discussed next, constitute about $45\%$, over 1.3 billion bases.
\item Ancient viral scraps account for about $8\%$, more than five times the entire protein-coding fraction.
\end{itemize}

For a specific example, human genomes contain over one million copies of a mobile element called \emph{Alu}, about 300 bases long. $10^{6}\times 300=3\times 10^{8}$ gives 300 million bases: this element alone occupies about a tenth of the genome. Calling the genome ``a book made of genes'' is therefore far from the full picture. It is more like a repeatedly photocopied working manual covered in sticky notes and stuffed with old advertisements, with the main instructions occupying only a few thin pages.
\end{qiaoliang}

\section{Mobile Paragraphs: Transposable Elements}

Copying a DNA segment requires at least a major accident such as unequal crossing over. Yet genomes contain even more restless residents with their own moving equipment: they cut themselves out, or copy themselves, and paste into another position. These are \textbf{transposable elements}, popularly ``jumping genes.''

\begin{zhengju}
In the 1940s, American scientist Barbara McClintock grew maize in Cold Spring Harbor's experimental fields. She focused on peculiar kernels with interlocking purple and white patches whose distributions could be inherited and predicted.

The standard explanation was that pigment genes were unstable and mutated readily. But under the microscope, McClintock saw overlooked chromosome details: a particular site repeatedly broke, and the breakage site \textbf{moved} between generations. She proposed a radical hypothesis: mobile ``controlling elements,'' named \emph{Ds} and \emph{Ac}, existed in genomes. When they jumped into a pigment gene, pigment production stopped and a white patch appeared; when they left, cells resumed making purple pigment. Patches were footprints of jumping events in cellular descendants.

Almost nobody believed this in the 1950s: textbooks described genes as fixed beads on a string, so how could they jump? McClintock temporarily stopped bothering to publish. Only when molecular biologists found mobile bacterial DNA in the 1960s and 1970s was her maize work rediscovered. In 1983, aged 81, she received an unshared Nobel Prize. Her method matters most: she did not discard strange observations but aligned three independent clues---patch distributions, chromosome breaks, and inheritance---to force a conclusion nobody then wanted to accept.
\end{zhengju}

We now know that transposable elements are astonishingly widespread. As calculated above, they have contributed nearly half the human genome. Most copies have accumulated disabling mutations and become immobile ``wreckage,'' but a few remain active. \emph{LINE-1} elements still jump to new positions in early human embryos. Insertions in unimportant regions may pass unnoticed; insertions in essential genes can cause disease. Some hemophilia and cancer cases originate in such events.

Evolution, however, never wastes old parts. Some incoming sequences happen to contain short gene-switching segments that the host recruits as regulatory elements. Selfish lodgers become punctuation in the instruction manual. Between disruption and recruitment, these self-moving paragraphs supply important variation.

\section{Borrowing a Neighbor's Page: Horizontal Transfer}

So far, changes have stayed ``within the family,'' passing from parents to offspring through vertical transmission. Life also has a sideways channel: \textbf{borrowing} genes directly from unrelated neighbors. This is \textbf{horizontal gene transfer}.

Bacteria are enthusiastic borrowers. They can transfer small circular DNA molecules, plasmids, through direct contact; pick up environmental DNA fragments from dead bacteria; or receive DNA carried as extra cargo by bacteriophages. Chapter 82's genetic engineering will revisit these three mechanisms. For now, remember that bacterial genomes are less private family property than a public lending library.

This directly affects us through antibiotic resistance. If a bacterium acquires a mutation producing an antibiotic-degrading enzyme, its offspring inherit resistance, the vertical story from Chapter 76. More troublingly, if the gene lies on a plasmid, the bacterium can \textbf{hand it directly} to an entirely different bacterial species. Hospital ``superbugs'' often arise not through patient accumulation of separate mutations but through repeated transfer and packaging of multiple resistance genes. That is why doctors emphasize adequate antibiotic doses and complete courses, preferring narrow-spectrum drugs when suitable: every misuse trains this gene-lending network through selection.

\textbf{Translator safety note:} Follow the clinician's current prescription, including dose and duration; do not select, stop, extend, or switch antibiotics yourself. ``Complete courses'' does not mean longer treatment always prevents resistance: clinicians should use an appropriate, evidence-based effective duration. See \href{https://www.cdc.gov/antibiotic-use/about/}{CDC patient guidance} and \href{https://www.nice.org.uk/guidance/ng15/evidence/full-guideline-pdf-252320799}{NICE antimicrobial stewardship guidance}.

Horizontal transfer is not confined to bacteria and sometimes crosses enormous gaps. This is our sweet potato's secret. Agrobacterium is a natural genetic engineer: when infecting plants, it inserts part of its DNA, called T-DNA, into plant genomes, making cells manufacture nutrients for it. Humans exploit this in genetic modification (Chapter 82), but the 2015 study found that Agrobacterium \textbf{had already done it for itself}. Every tested cultivated sweet-potato variety contained its T-DNA fragments, while wild relatives did not. This strongly suggests that borrowed bacterial genes contributed to the crucial trait of enlarged storage roots. The sweet potatoes humans have selected and enjoyed for millennia may themselves be products of natural horizontal transfer.

An important qualification: commonplace in bacteria, such events are very rare in animals and plants. Multicellular organisms have a ``wall'' separating reproductive and body cells (Chapter 75). Foreign DNA entering a toe cell cannot reach sperm or eggs and cannot be inherited. Sweet potatoes were fortunate that the insertion occurred in material capable of flowering, fruiting, and vegetative propagation through vine cuttings.

\section{Viral Fossils: Endogenous Sequences}

Horizontal transfer has a special courier: viruses.

Retroviruses, including HIV, have RNA genomes. After entering a cell, they reverse-transcribe RNA into DNA and \textbf{stitch} that DNA into host chromosomes. This reproductive strategy ensures that whenever the cell copies its DNA, it also copies the virus.

Imagine a rare but inevitable accident: insertion occurs in a germ cell that develops into a new individual. The viral ``advertisement'' then appears in every cell and passes to all descendants. Such remnants are \textbf{endogenous retroviral sequences}, which, as noted above, account for about 8\% of the human genome. Your genome therefore contains layered ``crime scenes'' from viral infections over tens of millions of years, recording ancestral epidemics like geological strata.

Most scraps are long since damaged and inactive. Evolution's tinkerer nevertheless reused them, this time with a remarkable result: your birth after growing in your mother's uterus may owe something to a viral gene.

A key placental structure forms when fetal cells \textbf{fuse} into a multinucleated syncytial layer, separating yet connecting maternal and fetal tissues so nutrients and oxygen can pass. The molecular ``zipper'' enabling fusion is a protein called syncytin, encoded by an ancient retroviral remnant. Viruses originally used it to penetrate cell membranes; mammals recruited it to build placentas. Later work found that different mammalian branches acquired syncytins from \textbf{different} ancient viruses, independently picking up the same kind of tool several times. A pathogenic virus's relic became a key component in the evolution of live birth.

In fairness, viruses do not only bring genes in; they sometimes carry host genes out in viral particles to the next host. Among microbes, viruses may be Earth's largest gene-moving service.

\section{Change the Instructions, Not the Parts}

Earlier sections concerned changes to sequence itself. Another evolutionary shortcut leaves the parts unchanged and modifies only their \textbf{operating instructions}.

Recall Chapters 74 and 75: nearby regulatory sequences determine which cells activate a gene, how strongly, and until what age. These sequences are DNA and can mutate. A small regulatory change can produce enormous morphological consequences because it alters not component quality but the \textbf{deployment plan}.

A classic case is the three-spined stickleback. Marine fish have large paired pelvic spines against predators. After the ice age ended about ten thousand years ago, many entered freshwater lakes, where numerous populations lost these spines. Genetic comparisons show \textbf{independent} losses in different lakes, often implicating the same gene, \emph{Pitx1}. Its coding sequence remains intact; the broken part is a nearby switch saying ``work in the pelvic-fin region.'' Why modify the switch rather than the gene? \emph{Pitx1} works elsewhere too, such as the brain and pectoral fins, so disabling it entirely causes collateral damage. Tearing out only the pelvic-fin instruction page minimizes side effects. Regulatory evolution thus permits \textbf{local} modification.

Human and chimpanzee protein-coding sequences are about 99\% similar. The other side of this familiar statistic is that much of our enormous difference concerns deployment plans, not parts lists. Let the same parts work for a few more cellular generations in the cerebral cortex, or switch off several days later in developing hands, and outcomes differ dramatically. Chapter 73's lactose tolerance is another regulatory example: variants telling the lactase gene to ``stay at work in adulthood'' were independently selected in different populations.

Follow this clue farther, and it reaches a deep evolutionary discovery: the master switches for eye formation in flies, mice, and humans descend directly from \textbf{the same gene}, called \emph{eyeless} in flies and \emph{Pax6} in vertebrates. Put mouse \emph{Pax6} into a fly, and it still produces eyes---fly-type eyes. Fly compound eyes and yours have no structural resemblance and evolved independently, yet the master switch activating eye construction in the right place existed over five hundred million years ago and has been used by all descendants ever since. Once evolution makes a useful switch, it refuses to rebuild it.

\section{A Tinkerer, Not a Designer}

We can now assemble the chapter's picture. In 1977, French biologist François Jacob, one of Chapter 70's operon-model proposers, wrote a famous essay whose title conveyed that evolution is a \textbf{tinkerer}.

An engineer draws plans, selects ideal parts, and manufactures from scratch. A tinkerer instead has a pile of old objects: spare duplicates, borrowed foreign components, and disabled viral wreckage. The method is hammering and patching, using old parts wherever they can serve a new purpose. Feathers may initially have insulated before switching careers to gliding and flight. Many transparent crystallins in eye lenses formerly did other jobs, such as enzymatic work, before being recruited to transmit light. Your vocal cords, jaw, and wrist bones all have predecessors in fish gill arches and fin rays, the homologous structures of Chapter 78.

This explains biology's famous inelegance: so many design flaws. The human retina is reversed, with photoreceptors \textbf{behind} blood vessels and nerves, forcing light through a layer of wiring and leaving a blind spot. The recurrent laryngeal nerve travels from the brain down beneath the aortic arch and back to the larynx, a four-meter detour in giraffes. No engineer would draw it this way, but a tinkerer must remodel inherited plans while keeping every intermediate viable. Evolution cannot demolish everything and restart.

Reusing old parts also implies that every present-day genome is an \textbf{archaeological accumulation}. At the bottom lies a billions-of-years-old core of replication and translation machinery shared by all life. Above it sit successive duplications, mobile elements, viral fragments, and rewritten switches. Reading this book reveals a chronicle of life; Chapter 81 explains how scientists actually read it.

\begin{wuqu}
``The human genome is 98\% junk DNA, all useless waste.'' Both ends of this statement are wrong. First, 98\% merely describes non-protein-coding sequence, which is far from synonymous with useless. Regulatory switches, functional RNAs, and chromosome structures such as centromeres and telomeres lie there, as do this chapter's syncytin and recruited transposon switches. But do not overcorrect to ``all junk DNA has a function.'' Genomes do contain much inert debris: disabled viral pages and mutated duplicate copies, mostly copied along without helping or hindering. The real picture is a storeroom containing tools, spare parts, and genuine rubbish. Evolution finds possibilities by rummaging through it. When someone claims all junk DNA is useful to support a startling conclusion or a perfectly designed genome, ask: what is the evidence?
\end{wuqu}

\begin{shiyan}
Use paper and pen to simulate copying, moving, and borrowing and experience accumulating variation. \textbf{Materials}: paper, a pen, a timer such as a phone, a sentence of about 30 characters copied from this book, and a family member or classmate. \textbf{Procedure}: (1) Show the sentence for 10 seconds, hide it, and ask the other person to reproduce it from memory. Compare the written version with the original and count ``mutations.'' (2) Play photocopier: the first person writes a sentence from memory, the second memorizes and rewrites that version, and so on through five people. Compare the final version with the original. (3) Simulate borrowing by inserting a sentence from another book into person three's version before passing it to person four. Do later writers treat it as part of the original? (4) Simulate duplication by writing one phrase twice in succession and observing whether both copies survive later transmission. \textbf{What to observe}: which changes persist most easily? After several rounds, does the inserted sentence still look foreign? \textbf{Think}: real DNA replication has proofreading and repair (Chapter 72); this game does not. What changes if each person can check a copy of the original and correct errors? \textbf{Safety reminder}: this paper-and-pen activity poses no risk. Avoid long sentences; 30 characters or fewer works best.
\end{shiyan}

\section{Where the Book Analogy Breaks Down}

Calling the genome a book has served us well, but now we must challenge the analogy. A book is inert; a genome is not. At least three differences can mislead. First, books do not decide which pages mutate readily. Genomic regions differ in change rates by hundreds or thousands of times because repair coverage and sequence chemistry affect where trouble occurs; changes are not evenly sprinkled like pepper. Second, borrowing book pages has no cost, whereas most real transfers and duplications offer no benefit or are harmful, and very few spread. Chapter 77's drift also reminds us that usefulness is not the only criterion for retention; luck matters greatly. Third, most easily forgotten, every major genomic event discussed here concerns accumulation in \textbf{populations} across generations, not one individual's lifetime. Your personal genome is largely fixed for life; the copying, borrowing, and remodeling belong to an inheritance chain extending hundreds of millions of years.

\section{Back to the Roasted Sweet Potato}

Break open the sweet potato again. It is sweet because heating breaks starch into smaller pieces and enzymes turn those into sugars, as Chapter 56 described. Its fluffy texture comes from starch granules filling the storage root. That ability to store nutrients in enlarged roots may well relate to foreign genes inserted by Agrobacterium hundreds of millions of years ago. Without laboratory procedures, those pages completed a textbook horizontal transfer: insertion, retention, spread through all cultivated varieties, and amplification by human selection into every field.

Next time you hear ``genetic modification,'' remember three things. First, genes have moved between species for hundreds of millions of years; humans merely learned to do it according to plans (Chapter 82). Second, an apparently foreign sequence can become essential to a species' new traits---our placentas provide evidence. Third, the questions for judging a technology are never simply ``natural or artificial?'' but ``What changed, what is the evidence, and what is the dose?''

\section{The Story Is Not Finished}

This chapter explored genomic hardware renovation: duplication, moving, borrowing, and altered switches. The biggest borrowing awaits next chapter: about two billion years ago, one cell swallowed an entire bacterium, did not digest it, and let it stay. That became the mitochondrion in each of your cells. Is borrowing a whole organelle the ultimate horizontal transfer or a new kind of cooperation? Chapter 80 enters cooperation, symbiosis, and multilevel selection. Chapter 81 then explains how we read these layered ancient texts, and Chapter 82 how tools from Agrobacterium and bacterial immunity let us rewrite the book ourselves.

\begin{tiaozhan}
What happens after a gene is duplicated? Population genetics identifies three main routes whose probabilities can be roughly estimated. With effective population size $N$, a neutral variant's chance of spreading to everyone is approximately its current frequency. A newly arisen copy has frequency $\frac{1}{2N}$, negligibly small. The first route, the fate of most duplicates, is quiet loss within a few generations or damaging mutations producing a pseudogene. Second, copies can \textbf{divide} the original work: a gene formerly active in two places acquires copies handling one each. This subfunctionalization can preserve both if each loses a different switch. The third, rarest, and most exciting route gives one copy a useful new function retained by selection: neofunctionalization, followed by red/green opsins and globin families. The proportions taking these routes remain an active question in evolutionary genomics. What is certain is that ``duplication to new function'' is no guaranteed script, but an occasional lottery win among a million duplications. Evolution's lottery budget consists of hundreds of millions of years and enormous populations.
\end{tiaozhan}

\section*{Try It Yourself}

\begin{sikaoti}
\item Red and green opsin genes are about 98\% similar and adjacent on the X chromosome. Draw their origin by duplication of one ancestral gene, showing two chromosomes, misalignment, and exchange. Which copy becomes ``free'' afterward? Why does the original resist accumulating large mutations?
\item McClintock's maize patches can be explained by genes jumping. Propose an alternative \textbf{without transposable elements}. Which observation, predictably inherited patches or moving chromosome-break sites, is hardest for your explanation to account for?
\item A patient has a superbug resistant to three antibiotics. Are its resistance genes more likely to have arisen through independent mutations or packaged plasmid borrowing? What epidemiological evidence would each hypothesis predict?
\item Syncytin came from an ancient virus. Someone concludes, ``Humans therefore evolved from viruses.'' Identify the level of the error and explain in a paragraph what viral sequence incorporation actually means.
\item In losing stickleback pelvic spines, evolution ``selected'' disruption of a \emph{Pitx1} regulatory switch rather than its coding sequence. Why does this route cost less? If a lake population instead lost spines through complete coding-sequence disruption, what other abnormalities might appear?
\item Imagine a future organism recruits an \emph{Alu} element from the human genome's storeroom as a useful switch. Is that element still junk? Organize your answer around who defines function and when.
\end{sikaoti}
